%%%%%%%%%%%%%%%%%%%%%%%%%%%% Document Setup %%%%%%%%%%%%%%%%%%
\documentclass[10pt]{article}
\usepackage[T2A,T1]{fontenc}
\usepackage[utf8]{inputenc}
\usepackage[russian,english]{babel}
\usepackage{CJKutf8}
\AtBeginDvi{\input{zhwinfonts}}
%\usepackage{polyglossia}
\makeatletter
\newlength{\bibhang}
\setlength{\bibhang}{1em}
\newlength{\bibsep}
 {\@listi \global\bibsep\itemsep \global\advance\bibsep by\parsep}
%\newenvironment{bibsection}%
 %       {\vspace{-\baselineskip}\begin{list}{}{%
   %    \setlength{\leftmargin}{\bibhang}%
     %  \setlength{\itemindent}{-\leftmargin}%
       %\setlength{\itemsep}{\bibsep}%
       %\setlength{\parsep}{\z@}%
        %\setlength{\partopsep}{0pt}%
        %\setlength{\topsep}{0pt}}}
        %{\end{list}\vspace{-.6\baselineskip}}
\makeatother
\reversemarginpar
\usepackage{geometry}\geometry{a4paper,total={170mm,257mm},left=1.18in,right=1.18in,top=0.79in,bottom=0.79in}
%\usepackage[paper=letterpaper,
           %\includefoot, % Uncomment to put page number above margin
   %         marginparwidth=0.6in, % Length of section titles
         %   marginparsep=0.05in, % Space between titles and text
      %      margin=1.18in, % 1 inch margins
            %includemp]{geometry}
%\usepackage{eurosans}
\setlength{\parindent}{0in}
\usepackage{enumerate} 
\usepackage{etaremune}
\usepackage{calc}
\usepackage{bbding}
\usepackage{paralist}
\usepackage{fancyhdr,lastpage}
\pagestyle{fancy}
%\pagestyle{empty} % Uncomment this to get rid of page numbers
\fancyhf{}\renewcommand{\headrulewidth}{0pt}
\fancyfootoffset{\marginparsep+\marginparwidth}
\newlength{\footpageshift}
\setlength{\footpageshift}
          {0.5\textwidth+1.\marginparsep+0.5\marginparwidth-2in}
\lfoot{\hspace{\footpageshift}%
       \parbox{4in}{\, \hfill %
                    \arabic{page} of \protect\pageref*{LastPage} % +LP
% \arabic{page} % -LP
                    \hfill \,}}
\usepackage{color,hyperref}
\usepackage{graphics}
\usepackage{wrapfig}
\usepackage{longtable}

\usepackage[usenames,dvipsnames,x11names,svgnames]{xcolor}
\definecolor{darkblue}{rgb}{0.0,0.0,0.3}
\definecolor{mygreen}{RGB}{44,204,148}
\definecolor{darkgreen}{RGB}{13,107,27}
\definecolor{lightgreen}{RGB}{0,170,37}
\definecolor{provinces}{RGB}{11,100,132}
\definecolor{mypurple}{RGB}{77,24,137}
\definecolor{haiyan}{RGB}{8,150,178}
\definecolor{anticred}{RGB}{165,62,11}
\hypersetup{colorlinks,breaklinks,
            linkcolor=darkblue,urlcolor=darkblue,
            anchorcolor=darkblue,citecolor=darkblue}
%%%%%%%%%%%%%%%%%%%%%%%% End Document Setup %%%%%%%%%%%%%%%%%%%
%%%%%%%%%%%%%%%%%%%%%%%%%%% Helper Commands %%%%%%%%%

\newcommand{\makeheading}[1]%
        {\hspace*{-\marginparsep minus \marginparwidth}%
         \begin{minipage}[t]{\textwidth+\marginparwidth+\marginparsep}%
                {\large \bfseries #1}\\[-0.15\baselineskip]%
                 \rule{\columnwidth}{1pt}%
         \end{minipage}}
\renewcommand{\section}[2]%
        {\pagebreak[2]\vspace{1.3\baselineskip}%
         \phantomsection\addcontentsline{toc}{section}{#1}%
         \hspace{0in}%
         \marginpar{
         \raggedright \scshape #1}#2}
% An itemize-style list with lots of space between items
\newenvironment{outerlist}[1][\enskip\textbullet]%
        {\begin{itemize}[#1]}{\end{itemize}%
         \vspace{-.6\baselineskip}}
% An environment IDENTICAL to outerlist that has better pre-list spacing
% when used as the first thing in a \section
\newenvironment{lonelist}[1][\enskip\textbullet]%
        {\vspace{-\baselineskip}\begin{list}{#1}{%
        \setlength{\partopsep}{0pt}%
        \setlength{\topsep}{0pt}}}
        {\end{list}\vspace{-.6\baselineskip}}

% An itemize-style list with little space between items
\newenvironment{innerlist}[1][\enskip\textbullet]%
        {\begin{compactitem}[#1]}{\end{compactitem}}
% An environment IDENTICAL to innerlist that has better pre-list spacing
% when used as the first thing in a \section
\newenvironment{loneinnerlist}[1][\enskip\textbullet]%
        {\vspace{-\baselineskip}\begin{compactitem}[#1]}
        {\end{compactitem}\vspace{-.6\baselineskip}}
% To add some paragraph space between lines.
% This also tells LaTeX to preferably break a page on one of these gaps
% if there is a needed pagebreak nearby.
\newcommand{\blankline}{\quad\pagebreak[2]}
% Uses hyperref to link DOI
\newcommand\doilink[1]{\href{http://dx.doi.org/#1}{#1}}
\newcommand\doi[1]{doi:\doilink{#1}}
% For \url{SOME_URL}, links SOME_URL to the url SOME_URL
\providecommand*\url[1]{\href{#1}{#1}}
% Same as above, but pretty-prints SOME_URL in teletype fixed-width font
\renewcommand*\url[1]{\href{#1}{\texttt{#1}}}

\usepackage{tikz}
\usetikzlibrary{chains}
\usepackage{pgfplots}
\usepackage{pgfmath}
\usetikzlibrary{matrix}
\usetikzlibrary{arrows,decorations.pathmorphing,backgrounds,positioning,fit,petri}
\usepgfmodule{decorations}
\usepgflibrary{decorations.shapes} 
\usetikzlibrary[decorations.shapes] 
\usetikzlibrary{mindmap}
\usepgflibrary{patterns}
\usetikzlibrary[patterns]
\usetikzlibrary{petri}
\usepgflibrary{plothandlers}
\usetikzlibrary{plothandlers}
\usepackage{pgfpages}
\usepgflibrary{arrows}
\usetikzlibrary{arrows}
\usetikzlibrary{trees}
\usetikzlibrary{topaths}
\usetikzlibrary{positioning}
\usetikzlibrary{fit}
\usepgflibrary{shapes.misc}
\usetikzlibrary{shapes.misc}
\usetikzlibrary{spy}
\usepgflibrary{shadings}
\usetikzlibrary{shadings}
\usepgfmodule{shapes}
\usepgfplotslibrary{external}
\usetikzlibrary{pgfplots.external}
\usepgfplotslibrary{colorbrewer}
\usepgfplotslibrary{patchplots}
\usetikzlibrary{plotmarks}
\usepackage{color} 
\usepackage{xcolor} 
%\pgfpagesuselayout{resize to}[a4paper,border shrink=20mm]
\pgfplotsset{colormap={cool}{rgb255(0cm)=(255,255,255); rgb255(1cm)=(0,128,255); rgb255(2cm)=(255,0,255)}}
  \pgfplotsset{colormap={violet}{rgb255=(25,25,122) color=(white) rgb255=(238,140,238)}}
  \pgfplotsset{
    colormap={bluered}{rgb255(0cm)=(0,0,180); rgb255(1cm)=(0,255,255); rgb255(2cm)=(100,255,0);
rgb255(3cm)=(255,255,0); rgb255(4cm)=(255,0,0); rgb255(5cm)=(128,0,0)}}
 \pgfplotsset{/pgfplots/colormap={winter}{rgb255=(0,0,255) rgb255=(0,255,128)}}
 \pgfplotsset{
    /pgfplots/colormap={spring}{rgb255=(255,0,255) rgb255=(255,255,0)}
}


\begin{document}
\selectlanguage{russian}

\begin{tikzpicture} %<- begin plot environment
\begin{axis}[
	title=Clusters from 15 experiment series,xlabel={\footnotesize{Humidity, \%}},
	ylabel={\tiny{Compressive strength, $C_{eq}$ (8h), MPa}},
	minor x tick num=5,minor y tick num=10,grid=major] %<- begin figure environment, parameters are given in square brackets
\addplot[ %<- parameters are given in square brackets
 	 scatter, %<- plot type - 'scatter'
    	only marks, %<- indicate that plot is made by dots, without connecting lines
    	point meta=explicit symbolic,
     	mark size=5pt,
    	scatter/classes={
		a={mark=+,DeepPink1}, 
		b={mark=o,DarkOrchid2}, 
		c={mark=x,DarkOrange3},
		d={mark=x,darkgreen},
		e={mark=Mercedes star,Green1},
		f={mark=star,SeaGreen2},
		g={mark=10-pointed star,Red2},
		h={mark=-,SpringGreen1},
		j={mark=10-pointed star,red},
		k={mark=asterisk,Purple3},
		l={mark=+,Magenta1},
		m={mark=-,OrangeRed1},
		n={mark=o,RoyalBlue1},
		o={mark=10-pointed star,Coral2},
		p={mark=-,Chartreuse2},
		r={mark=pentagon,DodgerBlue1} %<- indicate various types of points (design, color, size)
	}]
table [meta=class] {
	x  	y  		class
	20.6  0.222 	a
	20.6  0.211 	a
	20.6  0.217 	a
	20.6  0.219	a	
	20.6  0.172 	a
	20.6  0.167 	a
	34.1  0.074 	b
	34.1  0.101 	b	
	34.1  0.092 	b
	34.1	0.088  	b	
	34.1	0.092	b
	34.1	0.075	b
	34.0	0.060	c	
	34.0	0.095	c
	34.0	0.090	c	 
	34.0	0.089	c
	34.0	0.056	c
	34.0	0.083	c
	23.4	0.274	d
	23.4	0.296	d
	23.4	0.293	d	
	23.4	0.251	d
	23.4	0.294	d
	23.4	0.295	d
	25.0	0.107	e
	25.0	0.092	e
	25.0	0.153	e	
	25.0	0.138	e
	25.0 0.127	e
	25.0	0.151	e
	48.2	0.068	f
	48.2	0.123	f	
	48.2	0.094	f
	48.2	0.113	f
	48.2	0.123	f
	48.2	0.089	f
	26.6	0.184	g	
	26.6	0.162	g	
	26.6	0.162	g
	26.6	0.220	g
	26.6	0.137	g
	26.6	0.262	g
	34.6	0.098	h
	34.6	0.081	h
	34.6	0.098	h
	34.6	0.077	h
	34.6	0.059	h
	34.6	0.085	h
	35.3	0.151	j
	35.3	0.121	j
	35.3	0.129	j
	35.3	0.103	j
	35.3	0.110	j
	35.3	0.131	j
	29.4	0.132	k
	29.4	0.152	k
	29.4	0.088	k
	29.4	0.118	k
	29.4	0.099	k
	29.4	0.115	k
	49.3	0.163	l
	49.3	0.119	l
	49.3	0.133	l
	49.3	0.111	l
	49.3	0.112	l
	49.3	0.082	l
	34.5	0.082	m
	34.5 0.100	m
	34.5	0.061	m	
	34.5	0.083	m	
	34.5	0.102	m
	34.5	0.050	m
	47.7 0.083	n
	47.7 0.205	n
	47.7 0.090	n
	47.7 0.057	n
	47.7 0.080	n
	47.7 0.117		n
	22.5 0.083	o
	22.5 0.205	o
	22.5 0.090	o
	22.5 0.057	o
	22.5 0.080	o
	22.5 0.117		o
	22.5 0.075	p
	22.5 0.079	p
	22.5 0.076	p
	22.5 0.081	p
	22.5 0.084	p
	22.5 0.068	p
	23.5 0.071	r
	23.5 0.170	r
	23.5 0.193	r
	23.5 0.182	r
	23.5 0.178	r
	23.5 0.146	r
};
\end{axis}
\end{tikzpicture}
\begin{tikzpicture}
\begin{axis}[title=Clusters from 15 experiment series,xlabel={\footnotesize{Density, g / cm$^{3}$ }},ylabel={\footnotesize{$C_{eq}$ (8h.), MPa}},minor x tick num=5,minor y tick num=10,grid=major]
\addplot[
 	 scatter,
    	only marks,
    	point meta=explicit symbolic,
     	mark size=5pt,
    	scatter/classes={
		a={mark=+,magenta}, 
		b={mark=o,lime}, 
		c={mark=x,violet},
		d={mark=x,yellow},
		e={mark=Mercedes star,YellowOrange},
		f={mark=star,Maroon},
		g={mark=10-pointed star,ProcessBlue},
		h={mark=-,green},
		j={mark=10-pointed star,teal},
		k={mark=asterisk,orange},
		l={mark=+,ProcessBlue},
		m={mark=-,Yellow},
		n={mark=o,Cyan},
		o={mark=10-pointed star,LightSlateBlue},
		p={mark=-,PineGreen},
		r={mark=pentagon,Olive}
	}]
table [meta=class] {
	x  	y  		class
	1.99  0.222 	a
	1.99  0.211 	a
	1.99  0.217 	a
	1.99  0.219	a	
	1.99  0.172 	a
	1.99  0.167 	a
	1.80  0.074 	b
	1.80  0.101 	b	
	1.80  0.092 	b
	1.80	0.088  	b	
	1.80	0.092	b
	1.80	0.075	b
	1.75	0.060	c	
	1.75	0.095	c
	1.75	0.090	c	 
	1.75	0.089	c
	1.75	0.056	c
	1.75	0.083	c
	1.91	0.274	d
	1.91	0.296	d
	1.91	0.293	d	
	1.91	0.251	d
	1.91	0.294	d
	1.91	0.295	d
	1.86	0.107	e
	1.86	0.092	e
	1.86	0.153	e	
	1.86	0.138	e
	1.86 0.127	e
	1.86	0.151	e
	1.73	0.068	f
	1.73	0.123	f	
	1.73	0.094	f
	1.73	0.113	f
	1.73	0.123	f
	1.73	0.089	f
	1.87	0.184	g	
	1.87	0.162	g	
	1.87	0.162	g
	1.87	0.220	g
	1.87	0.137	g
	1.87	0.262	g
	1.77	0.098	h
	1.77	0.081	h
	1.77	0.098	h
	1.77	0.077	h
	1.77	0.059	h
	1.77	0.085	h
	1.61	0.151	j
	1.61	0.121	j
	1.61	0.129	j
	1.61	0.103	j
	1.61	0.110	j
	1.61	0.131	j
	1.85	0.132	k
	1.85	0.152	k
	1.85	0.088	k
	1.85	0.118	k
	1.85	0.099	k
	1.85	0.115	k
	1.65	0.163	l
	1.65	0.119	l
	1.65	0.133	l
	1.65	0.111	l
	1.65	0.112	l
	1.65	0.082	l
	1.69	0.082	m
	1.69 0.100	m
	1.69	0.061	m	
	1.69	0.083	m	
	1.69	0.102	m
	1.69	0.050	m
	1.51 0.083	n
	1.51 0.205	n
	1.51 0.090	n
	1.51 0.057	n
	1.51 0.080	n
	1.51 0.117		n
	1.51 0.083	o
	1.51 0.205	o
	1.51 0.090	o
	1.51 0.057	o
	1.51 0.080	o
	1.51 0.117		o
	1.89 0.075	p
	1.89 0.079	p
	1.89 0.076	p
	1.89 0.081	p
	1.89 0.084	p
	1.89 0.068	p
	1.60 0.071	r
	1.60 0.170	r
	1.60 0.193	r
	1.60 0.182	r
	1.60 0.178	r
	1.60 0.146	r
};
\end{axis}
\end{tikzpicture}
%	\begin{flushleft}
%		\selectlanguage{russian}\textbf{Рис. 3. График рассеяния величины восьмичасового эквивалентного сцепления грунтов с изменением влажности (слева) и плотности (справа) образцов грунтов}\\\vspace{6pt}
%		\textbf{Fig. 3. Scatter plot of the variations of the 8-hours equivalent cohesion of soil samples under changed humidity (left) and density (right) conditions}
%	\end{flushleft}
\end{document}
